% Part: proof-theory
% Chapter: proof-search
% Section: tableaux

\documentclass[../../../include/open-logic-section]{subfiles}

\begin{document}

\olfileid{pt}{ps}{tab}

\olsection{分析表}

分析表系统是与相继式演算密切相关的推导系统，特别适合于证明搜索，无论手动进行还是由软件实现。在分析表中，推导是一棵公式树，而非相继式树。但与自然演绎不同，这里的规则对应于相继式演算的规则。分析表推导本质上是对模型的显式搜索。因此它们有时被称为\emph{语义}分析表，或者也叫\emph{真值树}。

\emph{带号}分析表是带号公式的树，带号公式可以是 $\sFmla{\True}{!A}$ 或 $\sFmla{\False}{!A}$ 的形式。树向下生长。它们以指定待推导内容的公式作为起点。例如，要推导相继式 $!A, !B \Sequent !C, !D$，它将以 $\sFmla{\True}{!A}, \sFmla{\True}{!B},
\sFmla{\False}{!C}, \sFmla{\False}{!D}$ 开始。使 $!A$ 和 $!B$ 为真且使 $!C$ 和 $!D$ 为假的结构，就表明 $!A, !B \Sequent !C, !D$ 不是有效的。

分析表的叶节点（最底端节点）可以通过分支扩展规则进行延伸，这些规则向同一分支添加额外的节点，或在下方添加两个兄弟节点从而将该分支分裂为两个新分支。规则可以应用到同一分支中位于它上方的任何带号公式。经典分析表系统 \Log{Tc} 的分支扩展规则在 \olref{tab:Tc} 中给出。

\subfile{rules-Tc}

如你所见，这些规则就是 \Log{G3c} 的规则倒转过来，而符号 $\True$ 和 $\False$ 分别表示一个公式是相继式左边还是右边的主公式。

如果一个分析表的分支同时包含 $\sFmla{\True}{!A}$ 和 $\sFmla{\False}{!A}$，则该分支是\emph{封闭}的。如果每个分支都封闭，我们就说整个分析表是封闭的。例如，下面是一个针对 $!A \lor (!B \land !C) \Sequent (!A \lor !B) \land
(!A \lor !C)$ 的分析表，即从 $\sFmla{\True}{!A \lor (!B \land
!C)}$ 和 $\sFmla{\False}{(!A \lor !B) \land (!A \lor !C)}$ 开始：

\begin{maxwidth}
  \begin{tableau}{}
  [\sFmla{\True}{\formula{A} \lor (\formula{B} \land \formula{C})},
    just=\TAss, checked
    [\sFmla{\False}{(\formula{A} \lor \formula{B}) \land
        (\formula{A} \lor \formula{C})}, just=\TAss, checked
      [\sFmla{\True}{\formula{A}}, just={\TRule{\True}{\lor}[1]}
        [\sFmla{\False}{\formula{A} \lor \formula{B}},
          just={\TRule{\False}{\land}[2]}, checked
          [\sFmla{\False}{\formula{A}},
            just={\TRule{\False}{\lor}[4]}
            [\sFmla{\False}{\formula{B}},
              just={\TRule{\False}{\lor}[4]}, close]
          ]
        ]
        [\sFmla{\False}{\formula{A} \lor \formula{C}},
          just={\TRule{\False}{\land}[2]}, checked
          [\sFmla{\False}{\formula{A}},
            just={\TRule{\False}{\lor}[4]}
            [\sFmla{\False}{\formula{C}},
              just={\TRule{\False}{\lor}[4]}, close]
          ]
        ]
      ]
      [\sFmla{\True}{\formula{B} \land \formula{C}},
        just={\TRule{\True}{\lor}[1]}, checked
        [\sFmla{\False}{\formula{A} \lor \formula{B}},
          just={\TRule{\False}{\land}[2]}, checked
          [\sFmla{\True}{\formula{B}},
            just={\TRule{\True}{\land}[3]}, move by=2
            [\sFmla{\True}{\formula{C}},
              just={\TRule{\True}{\land}[3]}
              [\sFmla{\False}{\formula{A}},
                just={\TRule{\False}{\lor}[4]}
                [\sFmla{\False}{\formula{B}},
                  just={\TRule{\False}{\lor}[4]},close
                ]
              ]
            ]
          ]
        ]
        [\sFmla{\False}{\formula{A} \lor \formula{C}},
          just={\TRule{\False}{\land}[2]}, checked
          [\sFmla{\True}{\formula{B}},
            just={\TRule{\True}{\land}[3]}, move by=2
              [\sFmla{\True}{\formula{C}},
                just={\TRule{\True}{\land}[3]}
                [\sFmla{\False}{\formula{A}},
                  just={\TRule{\False}{\lor}[4]}
                  [\sFmla{\False}{\formula{C}},
                  just={\TRule{\False}{\lor}[4]},close
                ]
              ]
            ]
          ]
        ]
      ]
    ]
  ]
\end{tableau}
\end{maxwidth}

对勾号表示相应的规则已在下方所有分支上应用于该公式。符号 $\otimes$ 表示一个分支是封闭的。规则后面的行号用于指示该规则应用到哪些带号公式（即所指行上且位于同一分支中的公式）。

在分析表中进行证明搜索是容易的。我们分阶段生成一个分析表。
在第 $0$ 阶段，只需将起始公式写在顶部。
在第 $n+1$ 阶段，对每个叶节点进行如下操作：找出第一个（最上方的）尚未勾除的带号公式。
对该公式应用分支扩展规则。
一旦该规则在包含该带号公式的每个分支中都已被应用，就将其勾除。
（上面的例子就是用这个算法生成的。）

形式为 $\sFmla{\True}{\lforall}$ 和 $\sFmla{\False}{\lexists}$ 的带号公式永远不会被勾除，因此必须特殊处理。
若分支上存在这类公式，我们必须确保 $\TRule{\True}{\forall}$ 和 $\TRule{\False}{\lexists}$ 规则最终对所有适用的项都被应用。
例如，可以在处理完最上方非此形式的带号公式之后，在每个阶段对每个分支上的所有此类公式应用这些规则，以此来保证这一点。
所使用的项~$t$，是第一个能由该分支上已出现的常项符号（若无常项符号出现，则用 $\Obj c_0$）以及起始公式中出现的函数符号构造出的项，且需满足相应的带号公式 $\sFmla{\True}{!B(t)}$ 或~$\sFmla{\False}{!B(t)}$ 尚未出现在该分支上。

如果这个证明搜索没有产生封闭的分析表，那么它要么包含一个有限的开分支，其中所有非原子公式都已被勾除，
要么搜索产生一棵无限且有限分叉的树，而该树必然包含一个无限分支。
就像在 \olref[cpl]{sec} 中一样，我们可以定义一个结构~$\Struct{M}$，使得若分支上出现 $\sFmla{\True}{!A}$，则 $\Sat{M}{!A}$；
若出现 $\sFmla{\False}{!A}$，则 $\Sat/{M}{!A}$。（取 $\Theta =
\Setabs{!A}{\sFmla{\True}{!A} \text{ on the branch}}$ 和 $\Xi =
\Setabs{!A}{\sFmla{\False}{!A} \text{ on the branch}}$。）

通过这种证明搜索方法生成的分析表，只需为每个节点指派一个相继式，即可轻松转换为 \Log{G3c} 证明。
如果起始公式对应一个相继式 $\Gamma \Sequent \Delta$，那么将每个起始公式标注为 $\Gamma \Sequent \Delta$。
如果通过对带号公式 $\sFmla{\True}{!A}$ 应用分支扩展规则来延伸分支，那么该叶节点标注为 $!A, \Pi \Sequent \Lambda$；
如果是对带号公式 $\sFmla{\False}{!A}$ 应用，则标注为 $\Pi \Sequent \Lambda, !A$。
在分析表应用 \TRule{\True}{\star} 规则之处，就对相继式应用 \LeftR{\star} 规则（以 $!A$ 为主公式）；
在应用 \TRule{\False}{\star} 规则之处，就应用 \RightR{\star} 规则。
该规则的前提就是添加到分析表中的相应节点上所标注的相继式。
一旦分析表因 $\sFmla{\True}{!A}$ 和 $\sFmla{\False}{!A}$ 而封闭，最终叶节点的标注将是形如 $!A, \Pi \Sequent \Lambda, !A$ 的相继式。
分析表中某些相邻的节点会被指派相同的相继式；删去重复的标注即可。
例如，上面的分析表可转换为下面的推导：

\begin{maxwidth}
$
\Axiom$!A \fCenter !A, !B$
\RightLabel{\RightR{\lor}}
\UnaryInf$!A \fCenter !A \lor !B$
\Axiom$!A \fCenter !A, !C$
\RightLabel{\RightR{\lor}}
\UnaryInf$!A \fCenter !A \lor !C$
\RightLabel{\RightR{\land}}
\BinaryInf$!A \fCenter (!A \lor !B) \land (!A \lor !C)$
\Axiom$!B, !C \fCenter !A, !B$
\RightLabel{\RightR{\lor}}
\UnaryInf$!B, !C \fCenter !A \lor !B$
\RightLabel{\LeftR{\land}}
\UnaryInf$!B \land !C \fCenter !A \lor !B$
\Axiom$!B, !C \fCenter !A, !C$
\RightLabel{\RightR{\lor}}
\UnaryInf$!B, !C \fCenter !A \lor !C$
\RightLabel{\LeftR{\land}}
\UnaryInf$!B \land !C \fCenter !A \lor !C$
\RightLabel{\RightR{\land}}
\BinaryInf$!B \land !C \fCenter (!A \lor !B) \land
(!A \lor !C)$
\RightLabel{\LeftR{\lor}}
\BinaryInf$!A \lor (!B \land !C) \fCenter (!A \lor !B) \land
(!A \lor !C)$
\DisplayProof$
\end{maxwidth}

\end{document}